% Generated draft: requires mathematical and source review.
\documentclass[11pt]{article}
\usepackage[margin=25mm]{geometry}
\usepackage{fontspec}
\setmainfont{Latin Modern Roman}
\usepackage{amsmath,amssymb,mathtools,unicode-math}
\setmathfont{Latin Modern Math}
\usepackage{enumitem,needspace,xurl,hyperref}
\setlength{\parindent}{0pt}
\setlength{\parskip}{5pt}
\setlength{\emergencystretch}{3em}
\newcommand{\N}{\mathbb{N}}
\newcommand{\Z}{\mathbb{Z}}
\newcommand{\Q}{\mathbb{Q}}
\newcommand{\R}{\mathbb{R}}
\newcommand{\C}{\mathbb{C}}
\newcommand{\F}{\mathbb{F}}
\newcommand{\A}{\mathbb{A}}
\newcommand{\PP}{\mathbb{P}}
\newcommand{\E}{\mathbb{E}}
\newcommand{\eps}{\varepsilon}
\newcommand{\dd}{\,\mathrm{d}}
\begin{document}
\section*{UnsolvedMath ID: 701762}
\textbf{Source problem code:} catalogue entry
\subsection{Definitions and mathematical statement}
Work in the usual first-order language of groups \(L_{\mathrm{grp}}=\{\cdot,{}^{-1},1\}\), with equality and without named parameters. For a group \(H\), define its elementary theory and universal theory by
\[
\operatorname{Th}(H)=\{\sigma:\sigma\text{ is an }L_{\mathrm{grp}}\text{-sentence and }H\models\sigma\},
\]
\[
\operatorname{Th}_{\forall}(H)=\{\forall x_1\cdots\forall x_n\,\varphi:\ n\geq 0,\ \varphi\text{ is quantifier-free with free variables among }x_1,\ldots,x_n,\ H\models\forall x_1\cdots\forall x_n\,\varphi\}.
\]
A theory is decidable if an algorithm halts on every input sentence and correctly determines membership in that theory. Decidability here concerns each fixed group; no uniform algorithm taking a group presentation as input is stipulated.

Let \(G\) be any finitely generated group with a direct product decomposition \(G=G_1\times G_2\). No nontriviality condition on the factors is stated in the source. The following are posed as well-known questions:
\begin{enumerate}
\item[(a)] Does every such decomposition satisfy
\[
\operatorname{Th}(G)\text{ is decidable}
\quad\Longleftrightarrow\quad
\operatorname{Th}(G_1)\text{ and }\operatorname{Th}(G_2)\text{ are both decidable}?
\]
\item[(b)] Does every such decomposition satisfy
\[
\operatorname{Th}_{\forall}(G)\text{ is decidable}
\quad\Longleftrightarrow\quad
\operatorname{Th}_{\forall}(G_1)\text{ and }\operatorname{Th}_{\forall}(G_2)\text{ are both decidable}?
\]
\end{enumerate}
The source states that an affirmative answer to (a) for almost soluble groups follows from the cited work of G. A. Noskov. Under the usual terminology, an almost soluble group is one having a soluble subgroup of finite index. A group \(S\) is soluble if its derived series, defined by \(S^{(0)}=S\) and \(S^{(k+1)}=[S^{(k)},S^{(k)}]\), reaches the trivial group after finitely many steps. The source does not explicitly specify the logical language or define almost soluble; the conventions above require review. The restricted affirmative assertion and the registry's open status are supplied source assessments, not independently verified conclusions.
\subsection{Short English statement}
For a finitely generated direct product \(G=G_1\times G_2\), ask separately whether decidability of all first-order group sentences, and decidability of universal group sentences, are each equivalent to the corresponding decidability for both factors. The source attributes these well-known questions to E. I. Timoshenko and reports an affirmative answer to (a) for almost soluble groups via Noskov.
\subsection{Sources}
\begin{itemize}
\item Local UnsolvedMath registry, permanent ID \(701762\), exact source code \texttt{catalogue entry}, title \emph{Kourovka Notebook Problem 20.104}. The supplied record labels the problem open; this label has not been independently verified. Record updated \texttt{2026-10-16T14:00:00Z}; upstream revision \texttt{d371ee917ac9647fbb06cd6fc529deec6f3dda82}.
\item \emph{The Kourovka Notebook (Unsolved Problems in Group Theory)}, problem 20.104, proposed by E. I. Timoshenko, 2022. The supplied background identifies the 20th issue (2022) and the 21st issue edition (2026), with identifier \texttt{arXiv:1401.0300v48}. The supplied collection description names E. I. Khukhro and V. D. Mazurov as editors.
\item G. A. Noskov, \emph{Izv. Akad. Nauk SSSR, Ser. Mat.}, \textbf{47}, no. 3 (1983), 498--517. The supplied statement cites this work for the affirmative answer to (a) for almost soluble groups. No article title or URL is supplied, and the cited implication has not been independently checked.
\end{itemize}
\end{document}
